\begin{center}\emph{General results on representations of numbers by forms and their connection with transformations.}\end{center}

\Article{166}

\emph{If the form $F$ contains the form $F'$, every number that can be represented by $F'$ can also be represented by $F$.}

Let the indeterminates of the forms $F,F'$ be $x,y;x',y'$ respectively. Suppose the number $M$ is represented by $F'$ on putting $x'=m$, $y'=n$, and that the form $F$ changes into $F'$ under the substitution

\[
x=\alpha x'+\beta y',\quad y=\gamma x'+\delta y'
\]

\SourceLeaf{147}

Then clearly, on putting

\[
x=\alpha m+\beta n,\quad y=\gamma m+\delta n
\]

$F$ takes the value $M$.

If $M$ can be represented in several ways by the form $F'$, for example also on putting $x'=m'$, $y'=n'$, several representations of $M$ by $F$ follow. For if both

\[
\alpha m+\beta n=\alpha m'+\beta n'\quad\text{and}\quad\gamma m+\delta n=\gamma m'+\delta n'
\]

we would have either $\alpha\delta-\beta\gamma=0$, making the determinant of the form $F=0$, contrary to hypothesis, or $m=m'$, $n=n'$. So $M$ can be represented in at least as many different ways by $F$ as by $F'$.

Thus, if both $F$ contains $F'$ and $F'$ contains $F$, \emph{that is}, if $F,F'$ are equivalent, and the number $M$ can be represented by either form, it can also be represented by the other, and in just as many different ways by one as by the other.

Finally, we note that in this case the greatest common divisor of the numbers $m,n$ equals the greatest common divisor of $\alpha m+\beta n$, $\gamma m+\delta n$. Let the former $=\Delta$, and choose the numbers $\mu,\nu$ so that $\mu m+\nu n=\Delta$. Then

\[
\begin{gathered}(\delta\mu-\gamma\nu)(\alpha m+\beta n)-(\beta\mu-\alpha\nu)(\gamma m+\delta n)\\=(\alpha\delta-\beta\gamma)(\mu m+\nu n)=\pm\Delta\end{gathered}
\]

So the greatest common divisor of $\alpha m+\beta n$, $\gamma m+\delta n$ divides $\Delta$. But $\Delta$ also divides it, since it clearly divides $\alpha m+\beta n$, $\gamma m+\delta n$. The greatest common divisor must therefore be $=\Delta.$ --- Thus, if $m,n$ are relatively prime, $\alpha m+\beta n.\quad\gamma m+\delta n$ are also relatively prime.

\Needspace{12\baselineskip}
\Article{167}

\textsc{Theorem.} \emph{Suppose the forms}

\[
axx+2bxy+cyy\quad .\ .\ .\ .\ .\ .\ (F)
\]

\[
a'x'x'+2b'x'y'+c'y'y'\quad .\ .\ .\ .\ .\ .\ (F')
\]

\emph{are equivalent, their determinant $=D$, and the latter changes into the former on putting}

\[
x'=\alpha x+\beta y,\quad y'=\gamma x+\delta y
\]

\SourceLeaf{148}

\emph{Suppose also that the number $M$ is represented by $F$ on putting $x=m$, $y=n$, and so by $F'$ on putting}

\[
x'=\alpha m+\beta n=m',\quad y'=\gamma m+\delta n=n'
\]

\emph{with $m$ relatively prime to $n$, and therefore also $m'$ to $n'$. Then the two representations belong either to the same value of the expression $\surd D(\operatorname{mod.}M)$ or to opposite values, according as the transformation of the form $F'$ into $F$ is proper or improper.}

\emph{Proof.} Choose the numbers $\mu,\nu$ so that $\mu m+\nu n=1$, and put

\[
\frac{\delta\mu-\gamma\nu}{\alpha\delta-\beta\gamma}=\mu',\quad \frac{-\beta\mu+\alpha\nu}{\alpha\delta-\beta\gamma}=\nu'
\]

These are integers because $\alpha\delta-\gamma\beta=\pm1$. Then

\[
\mu'm'+\nu'n'=1.\quad\text{(See the end of the preceding article.)}
\]

Further, put

\[
\begin{gathered}\mu(bm+cn)-\nu(am+bn)=V,\\\mu'(b'm'+c'n')-\nu'(a'm'+b'n')=V'\end{gathered}
\]

Then $V,V'$ are values of the expression $\surd D(\operatorname{mod.}M)$ to which the first and second representations belong. If in $V'$ we substitute the values of $\mu',\nu',m',n'$, and in $V$ substitute

\[
\begin{array}{lll}\text{for}&a,&a'\alpha\alpha+2b'\alpha\gamma+c'\gamma\gamma\\\text{for}&b,&a'a\beta+b'(\alpha\delta+\beta\gamma)+c'\gamma\delta\\\text{for}&c,&a'\beta\beta+2b'\beta\delta+c'\delta\delta\end{array}
\]

expansion gives $V=V'(\alpha\delta-\beta\gamma).$

So either $V=V'$ or $V=-V'$, according as $\alpha\delta-\beta\gamma=+1$ or $=-1$. \emph{That is}, the representations belong to the same value of the expression $\surd D(\operatorname{mod.}M)$ or to opposite values, according as the transformation of the form $F'$ into $F$ is proper or improper. $Q.\ E.\ D.$

Thus, if several representations of the number $M$ by the form $(a,b,c)$, using relatively prime values of the indeterminates $x,y$, belong to \emph{different} values of the expression $\surd D(\operatorname{mod.}M)$, the corresponding representations by the form $(a',b',c')$ belong to those same respective values. If there is no representation of the number $M$ by a form belonging to a particular value, there is also none belonging to this value by an equivalent form.

\SourceLeaf{149}

\Article{168}

\textsc{Theorem.} \emph{If the number $M$ is represented by the form $axx+2bxy+cyy$ by assigning to $x,y$ relatively prime values $m,n$, and the value of the expression $\surd D(\operatorname{mod.}M)$ to which this representation belongs is $N$, then the forms $(a,b,c)$, $(M,N,\frac{NN-D}{M})$ are properly equivalent.}

\emph{Proof.} Article $155$ shows that integers $\mu,\nu$ can be found so that

\[
m\mu+n\nu=1,\quad \mu(bm+cn)-\nu(am+bn)=N
\]

Then the form $(a,b,c)$, under the substitution $x=mx'-\nu y'$, $y=nx'+\mu y'$, which is clearly proper, changes into a form whose determinant $=$ $D(m\mu+n\nu)^2$, \emph{that is}, $=D$, and so into an equivalent form. If we write this form $=(M',N',\frac{N'N'-D}{M'})$, we have

\[
\begin{gathered}M'=amm+2bmn+cnn=M,\\N'=-m\nu a+(m\mu-n\nu)b+n\mu c=N\end{gathered}
\]

\begingroup\rightskip=0pt plus 3em\relax
So the form into which $(a,b,c)$ changes under that transformation is $(M,N,\frac{NN-D}{M})$. $Q.\ E.\ D.$
\par\endgroup

The equations

\[
m\mu+n\nu=1,\quad \mu(mb+nc)-\nu(ma+nb)=N
\]

also give

\[
\begin{gathered}\mu=\frac{nN+ma+nb}{amm+2bmn+cnn}=\frac{nN+ma+nb}{M},\\\nu=\frac{mb+nc-mN}{M}\end{gathered}
\]

so these numbers are integers.

We must also note that this proposition does not apply if $M=0$, since the term $\frac{NN-D}{M}$ is then indeterminate\HistoricalNote{*)}{In this case, if we wish to extend our wording to it, saying that $N$ is a value of the expression $\surd D(\operatorname{mod.}M)$, or that $NN\equiv D(\operatorname{mod.}M)$, means that $NN-D$ is a multiple of $M$, and so $=0$.}.

\Article{169}

If we have several representations of the number $M$ by $(a,b,c)$ belonging to the same value of the expression $\surd D(\operatorname{mod.}M)$, $N$ (always assuming the values of $x,y$ to be relatively prime), several proper transformations of the form $(a,b,c)..(F)$ into $(M.\ N,\frac{NN-D}{M})..(G)$ can also be derived from them. For if such a representation also comes from the values $x=m'$, $y=n'$, then $(F)$ also changes under the substitution

\SourceLeaf{150}

\[
\begin{gathered}x=m'x'+\frac{m'N-m'b-n'c}{M}y',\\y=n'x'+\frac{n'N+m'a+n'b}{M}y'\end{gathered}
\]

into $(G)$. Conversely, every proper transformation of the form $(F)$ into $(G)$ gives a representation of the number $M$ by the form $(F)$ belonging to the value $N$. For if $(F)$ changes into $(G)$ on putting $x=mx'-\nu y'$, $y=nx'+\mu y'$, then $M$ is represented by $(F)$ on putting $x=m$, $y=n$. Since here $m\mu+n\nu=1$, the value of the expression $\surd D(\operatorname{mod.}M)$ to which the representation belongs is $\mu(bm+cn)-\nu(am+bn)$, \emph{that is}, $N$. Several different proper transformations give just as many different representations belonging to $N$\HistoricalNote{*)}{If two different proper transformations are assumed to give the same representation, they must have the following form:
\[
\begin{aligned}
1)\quad &x=mx'-\nu y',\quad y=nx'+\mu y';\\
2)\quad &x=mx'-\nu'y',\quad y=nx'+\mu'x'.
\end{aligned}
\]
But the two equations
\[
\begin{gathered}
m\mu+n\nu=m\mu'+n\nu',\\
\mu(mb+nc)-\nu(ma+nb)=\mu'(mb+nc)-\nu'(ma+nb),
\end{gathered}
\]
easily give either $M=0$ or $\mu=\mu'$, $\nu=\nu'$. We have already excluded $M=0$.}. --- It follows easily that once all proper transformations of the form $(F)$ into $(G)$ are known, they give all representations of $M$ by $(F)$ belonging to the value $N$. So the problem of finding representations of a given number by a given form, in which the indeterminates have relatively prime values, has been reduced to finding all proper transformations of that form into a given equivalent form.

Applying the results of Article $162$, we easily conclude: if one representation of the number $M$ by the form $(F)$ belonging to the value $N$ is $x=\alpha$, $y=\gamma$, then the general formula including all representations of the same number by the form $(F)$ belonging to the value $N$ is

\[
x=\frac{\alpha t-(\alpha b+\gamma c)u}{m},\quad y=\frac{\gamma t+(\alpha a+\gamma b)u}{m}
\]

where $m$ is the greatest common divisor of the numbers $a,2b,c$, and $t,u$ stand for all numbers satisfying the equation $tt-Duu=mm$.

\Article{170}

If the form $(a,b,c)$ is equivalent to an ambiguous form, and so is both properly and improperly equivalent to the form $(M,N,\frac{NN-D}{M})$, or properly equivalent both to the form $(M,N,\frac{NN-D}{M})$ and to $(M,-N,\frac{NN-D}{M})$, there are representations of the number $M$ by the form \SourceLeaf{151}$(F)$ belonging both to the value $N$ and to the value $-N$. Conversely, if several representations of the number $M$ by the same form $(F)$ belong to \emph{opposite} values of the expression $\surd D(\operatorname{mod.}M)$, $N,-N$, then the form $(F)$ is both properly and improperly equivalent to the form $(G)$, and an ambiguous form equivalent to $(F)$ can be found.

\Needspace{7\baselineskip}
Let these general results on representations suffice here. We will discuss representations in which the indeterminates do not have relatively prime values below. As for other properties, forms with negative determinant need a completely different treatment from forms with positive determinant, so we will now consider them separately. We begin with the former, as the easier ones.
